\documentclass[11pt,a4paper]{article}

\usepackage[margin=2.5cm]{geometry}
\usepackage[T1]{fontenc}
\usepackage[utf8]{inputenc}
\usepackage{lmodern}
\usepackage{hyperref}
\usepackage{graphicx}
\usepackage{enumitem}
\usepackage{titlesec}

\hypersetup{
  colorlinks=true,
  linkcolor=blue,
  urlcolor=blue,
  citecolor=blue,
  pdftitle={Computer Science Portfolio Report}
}

\title{Computer Science Portfolio Report}
\author{Aiwyn Abraham Anish}
\date{\today}

\begin{document}

\maketitle

\section{Introduction}

This report describes the design, implementation, and content of my
computer science portfolio. The portfolio was created to present my
technical skills, academic development, and practical programming work
for working student opportunities in software engineering.

The live portfolio website is available at:

\begin{center}
\url{https://aiwyn-007.github.io/Portfolio/}
\end{center}

The corresponding GitHub repository is available at:

\begin{center}
\url{https://github.com/Aiwyn-007/Portfolio}
\end{center}

The portfolio combines a responsive personal website, technical project
materials, a LaTeX-created curriculum vitae, and this report. GitHub
Pages was selected because it provides a simple way to publish a static
website directly from a GitHub repository.

\section{Portfolio Website Structure}

The portfolio is implemented as a single-page website with anchor-based
navigation. The main sections are:

\begin{itemize}
    \item \textbf{Introduction section:} Introduces my professional goal as a
    future software engineer and provides quick navigation actions.
    \item \textbf{About section:} Summarises my interests also states my approach to
    learning and focus on efficient and user-centred software.
    \item \textbf{Skills section:} Presents programming languages, web
    fundamentals, development tools, and professional skills.
    \item \textbf{Education section:} States my BEng Software Engineering
    studies and expected graduation date.
    \item \textbf{Portfolio section:} Links to selected technical exercises
    and explains the technologies and outcomes of each project.
    \item \textbf{Contact section:} Provides a professional contact method.
\end{itemize}

The website is designed to be responsive. CSS grid layouts, flexible
navigation, and media queries allow the content to remain readable on
desktop and mobile screens.

\section{GitHub Repository Structure}

The repository is organised to separate website files, documents, and
technical exercises clearly:

\begin{verbatim}
portfolio/
├── index.html
├── README.md
├── assets/
├── cv/
└── report/
\end{verbatim}

The \texttt{index.html} file is the entry point for the website.
The \texttt{assets} directory stores reusable documents and images,
including the PDF version of my CV. The \texttt{projects} directory
contains individual technical exercises. Each project includes source
code and a README file explaining its purpose, setup, technologies, and
key learning outcomes. The \texttt{cv} directory contains both the
LaTeX source and compiled PDF of my CV. Finally, the \texttt{report}
directory contains the source and PDF of this report.

This structure improves maintainability because related files are grouped
together. It also helps a hiring manager quickly locate the live website,
project source code, CV, and documentation.

\section{Design Decisions}

A dark red-and-black visual theme was selected to create a modern
and consistent appearance.Red is used as an accent colour
for calls to action,headings and interactive elements while dark
backgrounds improve contrast and focus attention on the content.

The website uses a minimal single-page layout rather than many separate
pages. This decision reduces navigation complexity and enables a hiring
manager to review the most important information quickly. Clear headings,
cards, whitespace, and concise text make the content scannable.

The project cards are designed to focus on evidence of technical ability.
For each project, the portfolio explains the problem, technologies used,
and a link to the relevant source code. This is more useful than listing
skills without practical examples.

\section{Tools and Technologies}

The following tools and technologies were used:

\begin{itemize}
    \item \textbf{HTML:} Used to define the semantic structure and
    content of the portfolio website.
    \item \textbf{CSS:} Used for layout, responsive behaviour, colours,
    typography, cards, buttons, and media queries.
    \item \textbf{JavaScript:} Available for future interactive features;
    the current website prioritises a lightweight static implementation.
    \item \textbf{Git and GitHub:} Used for version control, repository
    hosting, and documentation.
    \item \textbf{GitHub Pages:} Used to deploy the static website from
    the repository.
    \item \textbf{LaTeX:} Used to create a structured and reproducible CV
    and project report.
\end{itemize}

GitHub Pages is appropriate for this project because the portfolio is a
static site made from HTML and CSS files. Publishing from the repository
also connects the live portfolio directly with its source code and
technical documentation.

\section{Strengths and Limitations}

A key strength of the portfolio is its clear visual hierarchy. The
website introduces my profile immediately, then presents skills,
education, projects, and contact information in a logical order. The
responsive layout makes the portfolio accessible across different screen
sizes. The GitHub repository also demonstrates organisation, version
control, and documentation practices.

The main current limitation is that, as an early-stage software
engineering student, the number and complexity of completed projects is
still growing.A portfolio is stronger when each listed skill is
supported by a visible project, test, demonstration, or code sample.
Another limitation is that the current website has limited interactivity,
which is an intentional trade-off to keep it lightweight and reliable.

\section{Future Improvements}

In future iterations,I would improve the portfolio by adding:

\begin{itemize}
    \item More completed projects with screenshots, live demos, and
    detailed technical documentation.
    \item A dedicated page for each project containing architecture,
    challenges, testing, and lessons learned.
    \item Automated testing and continuous integration for programming
    exercises.
    \item Accessibility improvements, including more detailed keyboard
    navigation testing and colour-contrast checks.
    \item A contact form implemented through a secure third-party service
    or backend solution.
    \item A downloadable CV button and links to professional profiles.
\end{itemize}

\section{Conclusion}

This portfolio provides a professional foundation for presenting my
technical development as a Software Engineering student.It combines a
responsive website,organised GitHub repository,LaTeX CV, and practical
technical exercises.This project will continue to evolve as I complete
new coursework and software projects.

\end{document}
